\section{Counting Rooms}
\label{sec:cses-counting-rooms}

\problemstatement{Given an $n\times m$ grid of floor (\texttt{.}) and wall
(\texttt{\#}) cells, count the number of rooms: maximal groups of floor
cells connected through shared edges (four directions).}

\begin{patternbox}
\emph{``Count connected regions in a grid''} is a \textbf{flood fill}:
whenever an unvisited floor cell is found, start a traversal that marks its
entire room, and increment a counter. Each new traversal is exactly one
room.
\end{patternbox}

\subsection*{Flood fill counts components}

Scan every cell. When you hit a floor cell not yet visited, it belongs to a
new room, so increment the room count and run a BFS or DFS that marks all
floor cells reachable from it (moving only up/down/left/right through
floor). Because each cell is marked the first time it is reached and never
revisited, the number of traversals launched equals the number of rooms.

\begin{ideabox}[One Launch per Component]
The visited grid guarantees each cell is consumed once. So the count of
``fresh starts'' -- unvisited floor cells that trigger a new flood -- is the
component count. This is the grid form of connected-component counting.
\end{ideabox}

\subsection*{Algorithm}

\begin{enumerate}
  \item For each cell $(i,j)$: if it is floor and unvisited, increment the
        room count and flood-fill from it.
  \item The flood marks every floor cell reachable through the four
        directions.
  \item Output the room count.
\end{enumerate}

\subsection*{Dry run}

\dryrun{$3\times3$ grid: row0 ``.\#.'', row1 ``.\#.'', row2 ``...''.}

\begin{itemize}
  \item Start at $(0,0)$: flood the left column and bottom row -- one
        room.
  \item $(0,2)$ connects down to $(1,2)$ then to the bottom row, already
        visited -- so it is the \emph{same} room via the bottom.
  \item Here everything floor is one connected room, so the answer is
        $1$.
\end{itemize}

\subsection*{C++ implementation}

\begin{lstlisting}
#include <bits/stdc++.h>
using namespace std;

int n, m;
vector<string> g;
vector<vector<bool>> vis;

void bfs(int si, int sj) {
    queue<pair<int,int>> q;
    q.push({si, sj}); vis[si][sj] = true;
    int dr[] = {-1,1,0,0}, dc[] = {0,0,-1,1};
    while (!q.empty()) {
        auto [r, c] = q.front(); q.pop();
        for (int k = 0; k < 4; k++) {
            int nr = r + dr[k], nc = c + dc[k];
            if (nr < 0 || nr >= n || nc < 0 || nc >= m) continue;
            if (g[nr][nc] == '#' || vis[nr][nc]) continue;
            vis[nr][nc] = true;
            q.push({nr, nc});
        }
    }
}

int main() {
    cin >> n >> m;
    g.resize(n);
    for (auto& row : g) cin >> row;
    vis.assign(n, vector<bool>(m, false));

    int rooms = 0;
    for (int i = 0; i < n; i++)
        for (int j = 0; j < m; j++)
            if (g[i][j] == '.' && !vis[i][j]) { rooms++; bfs(i, j); }

    cout << rooms << "\n";
    return 0;
}
\end{lstlisting}

\begin{complexitybox}
\textbf{Time Complexity.} $O(nm)$ -- each cell is visited once.
\textbf{Space Complexity.} $O(nm)$ for the visited grid and BFS queue.
\end{complexitybox}

\begin{takeawaybox}
Counting rooms is connected-component counting on an implicit grid graph:
one flood fill per unvisited region. This is the simplest graph traversal
and the template for every grid problem that follows.
\end{takeawaybox}
